\documentclass[10pt,twocolumn]{article}
\usepackage[margin=0.72in]{geometry}
\usepackage{amsmath,amssymb,amsthm,booktabs,array}
\usepackage[T1]{fontenc}
\usepackage{lmodern}
\usepackage[hidelinks,breaklinks]{hyperref}
\usepackage{microtype}
\usepackage{graphicx}
\usepackage{balance}
\graphicspath{{../../.worktree/logp/figures/}}
\usepackage[section]{placeins}
\setlength{\columnsep}{0.25in}
\setlength{\emergencystretch}{1em}
\raggedbottom
\newcolumntype{P}[1]{>{\raggedright\arraybackslash}p{#1}}
\newtheorem{proposition}{Proposition}
\newtheorem{theorem}[proposition]{Theorem}
\newcommand{\src}[1]{\nolinkurl{#1}}
\newcommand{\draftslot}[2]{%
  \typeout{DRAFT-SLOT: #1}%
  \par\smallskip{\raggedright\noindent\textbf{[PLACEHOLDER: #1]}\par}\nopagebreak[4]
  \noindent{\small #2}\par\smallskip}
\newcommand{\placeholderfigure}[3]{%
  \typeout{DRAFT-SLOT: #1}%
  \begin{figure}[t]\centering
  \fbox{\begin{minipage}{0.91\columnwidth}\raggedright
  \textbf{[PLACEHOLDER: #1]}\par\smallskip
  \small #2\par\medskip\textit{No measurements supplied.}
  \end{minipage}}
  \caption{#3}\label{fig:#1}\end{figure}}
\title{Fair Service, Local Execution:\\Usage-Fair Delegation Locks}
\author{Working paper draft}
\date{}
\begin{document}
\maketitle

\begin{abstract}
Shared locks schedule a scarce resource: serialized execution on protected
state. Usage-aware locks such as CFL distribute that resource by scheduling
lock owners. Yet selecting an owner also selects the core that executes its
work, potentially moving protected state between cores. Which costs belong
to fair service allocation, and which follow from this execution choice?
We develop a framework that separates the service contract from communication
and execution costs, then connects them through the executor sequence.
The analysis motivates usage-fair delegation: schedule requests by accumulated
service while retaining execution on a combiner. FC-PQ provides a concrete
design, and a compact cost model identifies when it can outperform CFL:
saved handoff costs must cover request scheduling and amortized combining
costs. Empirical results and implementation-refinement evidence remain
explicitly unfilled.
\end{abstract}

\section{Introduction}
Shared indexes, queues, and other mutable structures expose a familiar
multicore bottleneck: several clients compete for work that must execute
serially. A lock protects that work, but its policy also determines how the
serialized capacity is divided. Equal turns are not equal service. If one
client's update takes three times as long as another's lookup, alternating
them gives the updater three-quarters of the protected execution time.
High aggregate throughput can therefore coexist with unequal service.
Figure~\ref{fig:serviceallocation} makes the distinction visible: operation
counts and accumulated service answer different questions.

\begin{figure}[ht]
\centering
\includegraphics[width=\columnwidth]{service-allocation.pdf}
\caption{Equal turns versus service-aware allocation. Width encodes
illustrative useful service, not measured wall time. The lower row balances
totals over the shown window, not every prefix. The operation mixes differ:
this is an allocation example, not a throughput comparison.}
\label{fig:serviceallocation}
\end{figure}

Usage-aware locks address this allocation problem by accounting for service,
rather than relying only on acquisition order. SCL exposes how lock usage
can subvert CPU scheduling goals and controls lock opportunity~\cite{scl}.
CFL treats lock occupation as a schedulable resource, combining usage
accounting, group support, and NUMA-aware owner selection~\cite{cfl}.
These designs establish that access to a lock is a resource to schedule,
not merely a sequence of turns.

For caller-executed usage-fair locks, however, improving owner selection leaves
an architectural constraint: the chosen owner also executes the next critical
section. Serving a different client on another core can therefore move the
protected state as well as change who receives service. CFL reduces movement
through NUMA-aware scheduling, but same-socket handoff is still different
from retaining one executor. This is not evidence that fairness necessarily
hurts throughput---CFL reports improvements. It instead raises a question:
\emph{which costs are required to distribute service fairly, and which arise
from tying service allocation to execution placement?}

We develop a framework that separates the service contract from communication
and execution costs, then connects them through the executor sequence. It
explains when a bounded service gap constrains consecutive caller execution.
This motivates a different execution arrangement: \emph{delegation}, where
a combiner executes requests on behalf of clients. A new client can receive
service without requiring execution to change cores. We use FC-PQ to make
this choice concrete and derive when its request, accounting, and batching costs
are smaller than the handoff costs it avoids. The result is a conditional
performance comparison with CFL and explicit requirements for the design,
not a claim that fair scheduling or delegation is free.

\draftslot{MOTIVATION-RESULT}{Insert one application or controlled workload
showing heterogeneous service and the allocation problem. Compare CFL and
FC-PQ using matched work and resources, with per-client service alongside
performance; no motivating measurement has yet been supplied.}

\section{An analysis framework}
\label{sec:storyframework}
We describe a lock through two sequences: the clients whose requests receive
service and the cores that execute those requests. A workload supplies
operations, arrivals, placement, and protected/request footprints. The
protocol determines publication, selection, execution, completion, and return.
Keeping these events distinct lets us ask two independent questions before
connecting their answers.

\subsection{Service contract}
\label{sec:storyfairness}
Our target is equal protected-service allocation among equal-weight clients,
not equal operation counts. Consider a fixed group selectable at every
decision, starting with equal cumulative-service credits that never reset.
If increments are nonnegative and each is at most $C_{\max}$, selecting a
minimum-credit client keeps the spread of credits at most $C_{\max}$ after
every operation. Serving a minimum can raise it by at most one largest
operation; the other credits do not change. Appendix~\ref{app:fairness}
gives the general proof. No communication cost or execution placement enters
this argument.

The contract has to name its service metric and eligible clients. A timer
including request handling does not automatically measure useful protected
work; a pending invocation is not necessarily visible to selection. Admission,
resubmission, and accounting therefore need explicit rules. This contract is
a requirement on a scheduler, not yet a guarantee about FC-PQ or a wall-clock
response bound.

\subsection{Execution cost}
Performance analysis takes the service policy and executor sequence as inputs,
without assuming fairness. We adapt LogP's separation of communication
latency, processor overhead, and communication rate to coherent shared
memory~\cite{logp,coherence}. Publication and notification may overlap useful
execution; protected-state movement and scheduler work can instead delay the
next operation. Cached polling is not one network message per load.

We account separately for useful work, request traffic, scheduling, state
movement, and batch administration. Dependencies determine which delays
overlap; resource demand can limit rate even when latency is hidden. Only
exposed, non-overlapping costs enter the serialized execution path. Completion
also differs from caller return when an executor continues combining.
Section~\ref{sec:storycost} applies this accounting; Appendix~\ref{app:costs}
gives its detailed mapping.

\subsection{Connecting service to execution}
A service contract restricts permissible service sequences. An execution
architecture determines which executor sequences can realize them. For two
continuously backlogged clients on distinct fixed cores, equal positive
service increments, no credit resets, and a gap bounded after each operation
limit consecutive services to one client. With caller execution, this also
limits consecutive execution on one core (Appendix~\ref{app:residency}).
This is a restricted mechanism, not an impossibility theorem for all fair locks.

The design implication is to avoid making each required change in service
recipient a required change in executor. This does not remove the cost of
fair scheduling. It identifies which movement can be avoided, which costs
remain, and what a candidate solution must make cheap enough to compensate.

\section{Design from the analysis}
\label{sec:storyidea}
The framework suggests placing service selection at a retained executor,
rather than using it to transfer execution to each selected client. A
delegation-style lock provides this arrangement: clients submit operations,
and a combiner executes them and returns their results. Different clients
can receive service while the protected state is accessed by the same core.
Figure~\ref{fig:execution} shows the architectural choice.

Flat combining and queue-based combining already exploit this execution
locality~\cite{fc,cc}. TCLocks makes delegation available through ordinary
locking interfaces~\cite{tclocks}, and SCL already suggests usage-fair scheduling
of delegated requests~\cite{scl}. Our question is what service contract this
arrangement supports and when its remaining costs are worth paying, not
whether fairness and delegation can be combined at all.

\begin{figure}[t]
\centering
\includegraphics[width=\columnwidth]{execution-placement.pdf}
\caption{Same equal-cost service order, different executor sequences.
Each column is a service step, not a different physical core. Teal symbols
denote the same protected working set. This is a schematic, not a CFL trace
or a cache-residency guarantee. Request/result communication remains.}
\label{fig:execution}
\end{figure}

\paragraph{FC-PQ as a concrete design.}
\label{sec:storylock}
FC-PQ places a usage-ordered priority queue behind flat combining. Clients
publish operations; the combiner selects by recorded usage, executes the
delegate, charges its service, and publishes the result. The queue is
sequential under combiner ownership. This existing design makes the framework's
separation concrete.

Three design responsibilities follow from the analysis. \emph{Accounting}
must charge the stated service, rather than an unexplained timer interval.
\emph{Admission} must expose pending requests and assign new or returning
clients a justified credit baseline. \emph{Batch control} must amortize
execution costs while limiting how long an executor caller continues after
its own result is ready. Service fairness alone does not solve return latency.

Under exclusive combiner access, switching clients within a batch introduces
no necessary protected-state transfer to another executing core. This does
not guarantee L1 residency: operations can evict one another's data, payloads
can be large, and later batches can use another combiner. The cost analysis
must decide whether the retained execution is worth the remaining overhead.

Our source audit also prevents an unconditional implementation guarantee.
Admission, visibility, and measurement obligations remain unresolved
(Appendix~\ref{app:implementation}); the configured limit of 64 counts queue
pops, not completed requests.

\section{When does delegation pay?}
\label{sec:storycost}
Compare the same completed operation mix, hardware budget, and warm useful
work. For sustained contention with no uncovered idle time, approximate the
mean cost per completion as
\begin{equation}
\begin{aligned}
 t_{\mathrm{CFL}} &\simeq c+h+m,\\
 t_{\mathrm{PQ}} &\simeq c+d+\frac{a}{b}.
\end{aligned}
\label{eq:storycosts}
\end{equation}
Here $c$ is the common warm operation cost. CFL pays exposed protocol and
policy cost $h$, plus actual protected-state movement $m$, per completion.
FC-PQ pays per-request communication, selection, and accounting cost $d$.
Its mean pass cost $a$ includes administration and any protected-state or
scheduler-state reacquisition when the combiner changes; $b$ is the actual
number of completions per pass, averaged as a ratio of totals.

Only non-overlapping costs on the serialized execution path are added.
CFL's $m$ reflects its actual locality policy, including same-socket handoffs
and local reuse; it is not a cross-socket charge imposed on every operation.
Likewise, $a$ includes delegation's remaining data movement rather than
assuming it disappears. Appendix~\ref{app:costs} connects these terms to the
full communication and execution model.

The model gives a direct comparison. FC-PQ has lower time per completion when
\begin{equation}
 \boxed{m+h>d+\frac{a}{b}.}
 \label{eq:storycrossover}
\end{equation}
In words, \emph{CFL's per-operation handoff and policy costs must exceed
FC-PQ's per-request costs plus its amortized pass cost}. When these costs are
stable over a regime and $m+h>d$, sufficient batch occupancy produces an
advantage. When $d\ge m+h$, increasing the batch alone cannot produce a strict
win in this model. The appendix gives the crossover explicitly.
Figure~\ref{fig:costcrossover} illustrates both the amortization opportunity
and the per-request cost floor that batching cannot remove.

\begin{figure}[ht]
\centering
\includegraphics[width=\columnwidth]{cost-crossover.pdf}
\caption{Schematic of Equation~\eqref{eq:storycrossover}, not measurements.
With fixed $a>0$ and $H=h+m>d$, the modeled costs tie at $b^*=a/(H-d)$;
strict advantage starts to its right. If $d\ge H$, no batch can create
a strict win. Costs and executor-change behavior must remain applicable
as occupancy changes.}
\label{fig:costcrossover}
\end{figure}

This comparison predicts a benefit from reusable shared state and enough
pending work to amortize a pass. It also predicts failures: sparse activity
leaves little to amortize; large payloads increase request cost; a large
priority queue can increase scheduler cost; and effective CFL locality can
leave little movement to save. If the serialized execution path is the
bottleneck, its lower cost translates into higher predicted throughput.
Otherwise publication, memory traffic, or another resource may limit both
locks, and that bottleneck must be modeled separately.

This is an analytical condition, not a measured speedup. Prediction requires
independently established costs and a model of occupancy and executor changes;
Section~\ref{sec:storyevaluation} describes the test.

\paragraph{A second comparison: what does fairness add?}
Outperforming CFL does not establish a low fairness overhead. Compare FC-PQ
with a policy-disabled variant retaining the same publication and batching
infrastructure. Bookkeeping and selection are direct additions; changed
operation order, occupancy, and locality are separate effects. This matched
comparison measures the policy's price, not the remaining advantage over CFL.

\section{Evaluation}
\label{sec:storyevaluation}
All experimental results remain unfilled. The evaluation distinguishes service
fairness, conditional advantage over CFL, incremental policy cost, and
application benefit.

\draftslot{EVAL-SETUP}{Specify revisions, machine/compiler, placement and total
core budget, workload provenance, repetitions, uncertainty, and instrumentation
cost.}

\subsection{Does the scheduler distribute the intended service?}
Heterogeneous operations test whether balanced accounting corresponds to
balanced useful service. Traces relate selectable membership and every
charge to completed work, comparing billed and useful-service spread
(Section~\ref{sec:storyfairness}). The initial test
uses equal credits and a fixed continuously selectable cohort with bounded
nonnegative increments and no resets. Separate arrival/reactivation tests
examine admission rather than borrowing that guarantee. Agreement cannot
establish unconditional correctness of current FC-PQ.

\draftslot{EVAL-SERVICE}{Supply cohort and accounting evidence, per-class
service, spread trajectories, and admission failures; identify which contract
premise each violation challenges.}

\subsection{When does FC-PQ outperform CFL?}
The principal comparison will use CFL's specified fair-policy mode, aligning
the scheduled entity and service metric with FC-PQ. Offered work, completed
operation mix, and hardware resources must be matched; per-class completions
and useful service prevent cheap operations from inflating throughput.
Varying footprint, load, and placement tests
Equation~\eqref{eq:storycrossover}, including CFL's NUMA locality.
Actual completions, not queue pops or budgets, determine batch occupancy.

We will calibrate costs independently and predict batch occupancy and executor
changes before held-out runs in the model's stated regime. Measured traces
will diagnose errors, not count as predictions (Appendix~\ref{app:costs}).

\placeholderfigure{EVAL-CFL}{Matched-work CFL and FC-PQ throughput across
footprint/load regimes, with useful-service fairness alongside. Mark predicted
crossovers, uncertainty, and the region where FC-PQ loses.}{Planned comparison
of conditional performance advantage and service allocation.}

\subsection{What is the incremental price of fairness?}
We will compare fair and policy-disabled versions of the same delegation
protocol; original FC versus FC-PQ is not this control. Separate controlled
measurements will distinguish direct accounting and selection costs from
changes in order, batching, locality, and completed mix
(Appendix~\ref{app:evidence}).

\draftslot{EVAL-COST}{Supply matched-control costs, independent calibration,
held-out prediction errors and crossover failures, distinguishing direct
components from indirect execution changes.}

\subsection{Does the application retain the benefit?}
An application comparison must preserve semantics and return to heterogeneous
service demand, counting adaptation and marshalling costs. Per-class return
distributions will separate executor-callers from waiting callers: completion
need not mean the executor can return immediately. Sparse activity, capacity
eviction, large request metadata, and tight return budgets deliberately probe
where locality or amortization ceases to help.

\draftslot{EVAL-APPLICATION}{Supply a semantically equivalent application,
delegate ownership/affinity restrictions, useful-work throughput and role-aware
return evidence, including losing regimes and non-lock bottlenecks.}


\section{Relationship to prior work}
\label{sec:storyrelated}
The service objective matters when comparing mechanisms: SCL's lock
opportunity includes idle opportunity~\cite{scl}, whereas our target is
protected service. CFL's usage/group policy and NUMA locality must be credited
under an aligned objective~\cite{cfl}. FC-PQ uses explicit delegates rather
than TCLocks' transparent locking interface~\cite{tclocks}.

SCL already suggests usage-fair scheduling of delegated requests. Our proposed
contribution is the contract-to-cost-to-design connection and conditional
CFL comparison, not that combination alone. The framework adapts established
combining and communication models~\cite{fc,cc,logp,coherence}.
\draftslot{NOVELTY-POSITION}{Complete comparison with prior fair delegation,
cohorting, locality/fairness analyses, and predictive lock models before fixing
the final novelty claim. The three reference papers do not establish absence
of earlier results.}

\section{Conclusion}
Analyzing service allocation separately from execution cost exposes a design
choice: change who receives service without requiring execution to move.
FC-PQ makes that choice concrete. Its advantage over CFL depends on saved
handoff costs covering request scheduling and amortized pass costs.
Establishing the implementation contract and testing that boundary remain
the next evidence obligations. The goal is an explained operating regime
for fair local execution, not fairness for free.

\balance

\begin{thebibliography}{9}
\bibitem{scl} Y. Patel, L. Yang, L. Arulraj, A. C. Arpaci-Dusseau,
R. H. Arpaci-Dusseau, and M. M. Swift. Avoiding Scheduler Subversion using
Scheduler-Cooperative Locks. EuroSys, 2020.
\url{https://doi.org/10.1145/3342195.3387521}.
\bibitem{cfl} J. Park and Y. I. Eom. Locks as a Resource: Fairly Scheduling
Lock Occupation with CFL. PPoPP, 2024.
\url{https://doi.org/10.1145/3627535.3638477}.
\bibitem{fc} D. Hendler, I. Incze, N. Shavit, and M. Tzafrir. Flat Combining
and the Synchronization-Parallelism Tradeoff. SPAA, 2010, pp. 355--364.
\url{https://doi.org/10.1145/1810479.1810540}.
\bibitem{cc} P. Fatourou and N. D. Kallimanis. Revisiting the Combining
Synchronization Technique. PPoPP, 2012.
\bibitem{tclocks} V. Gupta, K. K. Dwivedi, Y. Kothari, Y. Pan, D. Zhou, and
S. Kashyap. Ship your Critical Section, Not Your Data: Enabling Transparent
Delegation with TCLocks. OSDI, 2023.
\url{https://www.usenix.org/conference/osdi23/presentation/gupta}.
\bibitem{logp} D. Culler et al. LogP: Towards a Realistic Model of Parallel
Computation. PPoPP, 1993. \url{https://doi.org/10.1145/155332.155333}.
\bibitem{coherence} S. Ramos and T. Hoefler. Modeling Communication in
Cache-Coherent SMP Systems---A Case-Study with Xeon Phi. HPDC, 2013.
\url{https://spcl.inf.ethz.ch/Publications/index.php?pub=162}.
\end{thebibliography}

\clearpage
\nobalance
\appendix
% Optional technical detail for the simpler paper-story companion.
\section{Cost accounting and the conditional comparison}
\label{app:costs}
For a finite trace, let $N>0$ be completions, $E>0$ passes, and $F$ executor
changes; count empty passes in $E$. Define ratios of totals
\begin{equation}
 \begin{gathered}
 b=\frac{N}{E},\qquad p=\frac{F}{E},\\
 a=A+p(M+K),\\
 \frac{a}{b}=\frac{EA+F(M+K)}{N}.
 \end{gathered}
 \label{eq:appendixmapping}
\end{equation}
Here $A$ is total exposed administration divided by $E$; $M$ and $K$ are
exposed protected-state and scheduler-state reacquisition totals divided by
$F$, respectively. When $F=0$, the change term is zero without defining
conditional means. Finite startup costs are separate. Warm operation cost
and $d$ are completion-weighted; unweighted averages of per-pass ratios are
incorrect. Administration includes work in empty passes, while operation
costs are attributed across completed work without discarding unsuccessful
protocol work.

This partition counts only exposed serialized-lane costs. Hidden communication
latency is not added again, although its resource demand can limit throughput.
A miss charged to service must not also enter movement or request cost.
The trace partition is exact accounting; replacing its terms by calibrated
parameters gives an approximation. With differing warm costs and uncovered
idle per completion, the fuller expressions are
\begin{equation}
 t_{\rm PQ}\simeq c_{\rm PQ}+d+a/b+I_{\rm PQ},\qquad
 t_{\rm CFL}\simeq c_{\rm CFL}+h+m+I_{\rm CFL}.
\end{equation}
A matched completed mix does not alone ensure equal warm costs: order,
footprint, and interference can matter. In the main text's common-$c$,
no-idle regime, cancellation gives Equation~\eqref{eq:storycrossover}.
For nonnegative costs, let $g=m+h-d$. If $g>0$, the equivalent threshold is
\begin{equation}
 b>\frac{a}{g}.
 \label{eq:appendixthreshold}
\end{equation}
If $g\le0$, no positive $b$ gives a strict advantage.
The mean $b$ is real, not necessarily integer or even at least one. Only
fixed positive integer completed batches permit the threshold
$b_{\min}=\lfloor a/g\rfloor+1$. Interpreting either threshold as a batch-size
prediction requires costs, including $p$, to remain applicable as $b$ changes.

CFL's $m$ reflects its actual executor sequence and topology, including
same-core reuse and same-node versus cross-node transfers. Its NUMA-aware
policy need not transfer state across sockets on every completion~\cite{cfl}.
Movement already exists in caller-executed locks; it is not all incremental
fairness overhead. The inequality orders modeled times, and orders throughput
only when the serialized lane is the bottleneck in a stable regime. Calibrated
predictions need actual protocol and hardware parameters. A guaranteed strict
machine ordering instead needs compatible bounds, for example an upper bound
on PQ time below a lower bound on CFL time. No numerical gain follows here.

\section{Independent service-fairness argument}
\label{app:fairness}
Fix a nonempty finite cohort $\mathcal A$, continuously selectable at every
selection boundary. Let $U_i$ be cumulative specified service, with initial
spread $D_0$. One selected client's service increases by a value in
$[0,C_{\max}]$ at each step; other credits remain unchanged, and none resets.
For fixed $\delta\ge0$, suppose selection satisfies
$U_i\le\min_{j\in\mathcal A}U_j+\delta$. Then, at every service boundary,
\begin{equation}
 \max_i U_i-\min_i U_i\le B_U:=\max\{D_0,\delta+C_{\max}\}.
 \label{eq:appendixspread}
\end{equation}
\begin{proof}
Initially the bound holds. If the old minimum is $u$, unchanged credits
remain in $[u,u+B_U]$. The selected credit cannot decrease and becomes at
most $u+\delta+C_{\max}\le u+B_U$. The new minimum cannot decrease;
therefore the spread stays at most $B_U$, completing induction.
\end{proof}
Exact-min selection and equal initial credits give $C_{\max}$. This proof
uses neither execution placement nor any performance parameter. Continuous
selectability is stronger than having a pending invocation; admission,
rejoining, and credit resets need separate contracts.

If useful service $U_i$ differs from billed credit $V_i=U_i+e_i$, assume
$\max_i e_i-\min_i e_i\le E_c$ at every selection. Choosing within $\delta$
of the minimum $V$ gives, for every $j$,
$U_i\le U_j+\delta+e_j-e_i\le U_j+\delta+E_c$.
Thus replace $\delta$ by $\delta+E_c$ above. Small per-operation errors do
not establish this cumulative pairwise bound. Without it, billed fairness
need not imply useful-service fairness. Nor does the service invariant
alone establish an elapsed-time deadline.

\section{Caller-execution residency}
\label{app:residency}
\begin{proposition}[Two-client residency constraint]
\label{prop:appendixresidency}
Two always-backlogged clients occupy distinct fixed cores. Each operation
executes on its caller's core, takes fixed protected service $c>0$, and adds
exactly $c$ to its persistent credit. Suppose $|U_1-U_2|\le B$ initially and
at every service boundary, with $B\ge c$. A same-client run has length at most
\begin{equation}
 q_{\max}=\left\lfloor\frac{2B}{c}\right\rfloor.
\end{equation}
If each executor change costs at least $M>0$ on the serialized lane,
non-overlapped with service or another change, then for every integer $n\ge1$,
\begin{equation}
 \begin{aligned}
 T(n)&\ge nc+M\left(\left\lceil\frac{n}{q_{\max}}\right\rceil-1\right),\\
 \liminf_{n\to\infty}\frac{T(n)}n&\ge c+\frac{M}{q_{\max}}.
 \end{aligned}
\end{equation}
\end{proposition}
\begin{proof}
At a run's start, its client's credit minus the other's is at least $-B$.
After $k$ services it has increased by $kc$ and cannot exceed $B$, so
$kc\le2B$. Consequently $n$ services require at least
$\lceil n/q_{\max}\rceil$ runs and one executor change between adjacent
runs. Add their stipulated non-overlapping costs; divide by $n$ and take
the lower limit.
\end{proof}
This is total modeled time, not incremental fairness overhead against an
unfair baseline or an empirical hardware theorem. A measured mean cannot
supply the stipulated lower bound $M$. Delegation removes the caller--executor
identity premise: different clients can be served on one core. It does not
ensure cache residency or remove request, metadata, and eviction costs.

Completion is not caller return. In a pass with actual $k\ge1$ completions,
an executor finishing its own request first may continue the pass. Only
certified subsequent-operation and remaining-work bounds justify
\begin{equation}
 R_{\rm tail}\le e_{\rm tail}+(k-1)\tau_{\max}.
\end{equation}
All post-completion work, including non-service attempts and suspension,
must be covered. A uniform envelope needs an enforced maximum $k$, never
mean $b$; service fairness supplies neither bound.

\section{Implementation obligations}
\label{app:implementation}
At audited baseline \texttt{56a02ab}, FC-PQ initializes zero-credit newcomers
with historical mean per-operation charge, not cumulative virtual service.
Its post-pop minimum clamp is a no-op under valid minimum ordering and
cannot promote an unpopped client. A reserved but unpublished announcement
can stall discovery; the pass limit counts 64 pops, not completions.
Billing surrounds more than benchmark useful work, and the combiner timer
resets its start during service, invalidating it as a full-pass measurement.
These facts prevent unconditional application of the abstract contracts.
\draftslot{IMPLEMENTATION-REFINEMENT}{Establish admission, persistent credits,
visibility progress, billing semantics, and correct measurement boundaries
in an identified implementation; supply correctness evidence before claiming
its guarantees. No production repair is part of this draft.}

\section{Evidence and reproduction}
\label{app:evidence}
The separate \src{performance.tex} and \src{fairness.tex} supply independent
analyses; \src{implementation-audit.tex} records source evidence.
\src{paper.tex} remains their driver. Build this companion through
\src{story-draft.tex}; \src{README.md} gives reproduction commands.
\src{check_model.py}'s existing eight groups do not verify the new residency
result or this CFL comparison, and prove no Rust or hardware guarantee.

Direct billing overhead needs fixed-selection replay. Direct selector costs
need identical selectable-state snapshots, pass boundaries, and controlled
metadata residency, without executing the selectors' chosen delegates.
End-to-end comparisons separately report changed order, mix, occupancy,
placement, and idle; equal source budgets do not ensure equal completed work.
\draftslot{ANALYTICAL-CHECKS}{Add finite residency checks, including nonintegral
$B/c$ and admissible initial gaps, and exact crossover checks. These supplement,
not replace, proofs and do not validate concrete concurrent execution.}
\draftslot{CALIBRATION-CONTRACT}{Calibrate protocol/hardware costs independently;
predict occupancy and executor changes before held-out runs. Supplying target
trace $b,p$ reconstructs that trace, not a held-out prediction. Report uncertainty
and failure regimes; no calibration or measured advantage is asserted.}

\end{document}
